\begin{textsample}{NEG Dim 5 – profiled\_gpt – Score -1.00 – profiled\_gpt/t001643\_gpt.txt}  \label{ex:f5_neg_039}
Well, I ’ll just say, to start with, that we ’re putting the recorder on now, and what we ’re going to be doing is talking together about our experiences of having, and being treated for, gynaecological cancers.
The idea is that this isn't just for us, but that the recording will be used for research connected with Cambridge University Press, and also for a charity called The Eve Appeal.
The Eve Appeal is very much about campaigning for \textbf{greater} awareness of gynaecological cancers among women, and in particular younger women.
There ’s this assumption that younger women are terribly open about sex and about their bodies, but in fact a lot of them don't really know the names of the various bits, or the proper terminology, and that can have quite serious consequences.
If you don't know what things are called, or you feel awkward about saying the words, you ’re less likely to go and get things checked, or even to notice when something ’s wrong.
So part of what we ’re doing here is not just telling our stories about cancer and treatment, but also, in a way, demonstrating that you can talk about these things plainly, and that it ’s important to do so.
It ties in with what The Eve Appeal is trying to get across : that knowing your own body, and having the language to describe it, is actually a matter of health, not just embarrassment or prudishness.
They put out this press release about all of that – about younger women, and about the lack of knowledge and the terminology – and it was quite striking.
I wondered, when you read The Eve Appeal ’s press release on this issue, what did you feel?

% matched lemmas: great
\end{textsample}
